%% =====================================================================
%%  T-06-02  The Two-Tellings Test
%%  -------------------------------------------------------------------
%%  One idea per file. Core is mandatory. Slots in canonical order.
%%  Max 700 words. No layout commands. No filler. New source material
%%  for this entry goes in sources/<BOOK>/T-06-02.tex -- never by
%%  rewriting this file (docs/RULES.md R0.1).
%% =====================================================================
%% @kind: topic
%% @id: T-06-02
%% @title: The Two-Tellings Test
%% @subtitle: Consistency across time is the cheapest honesty check
%% @chapter: c06
%% @order: 20
%% @status: stable
%% @sources: B1
%% @updated: 2026-09-19
%% @allow-rewrite: slot migration to the seven-question format (docs/RULES.md section 2)

\topic{T-06-02}{The Two-Tellings Test}{Consistency across time is the cheapest honesty check}


\begin{core}
Let a person tell the same story twice, at different times, and compare the details. A true account is stable in its specifics and varies in its framing; an invented one varies in its specifics because the specifics are being regenerated to suit whatever the current purpose is.
\end{core}

\begin{mechanism}
A remembered event is retrieved as a whole, so its concrete details --- who was there, what order things happened in, what was said --- come back the same each time even when emphasis and interpretation change. A constructed account is assembled on demand, and assembly is driven by present need; the details shift because they were never stored, only chosen. Habitual construction makes this worse rather than better, because a person who routinely reshapes stories to fit the moment loses track of what was said before. The test is cheap because it requires no confrontation and no expertise --- only memory and a second occasion.
\end{mechanism}

\begin{application}
  \item Ask the same question twice, separated by time, in different contexts.
  \item Record the specifics the first time. Do not signal that you are recording them.
  \item Compare specifics, not impressions. Framing differences are ordinary; factual differences are not.
  \item Where the account has shifted, ask about the specific rather than the story, and observe whether the response is recall or reconstruction.
\end{application}

\begin{feedback}
  \item Names, dates, amounts or sequences differ between tellings.
  \item The story improves: each version is more favourable to the teller than the last.
  \item Detail density changes --- one version is rich in irrelevant specifics, another is vague where the first was precise.
  \item The moral of the story changes with the audience.
  \item Challenges to a specific are met with irritation rather than with recall.
\end{feedback}

\begin{failure}
It requires repeated contact, so it is useless in one-off dealings --- which is where most losses occur. It also fails in two directions: an honest person under stress, or telling a story they were not present for, produces exactly the same drift; and a careful liar who keeps a consistent script passes it easily. It is a screening test with modest accuracy, not a determination, and treating it as the latter produces false accusations against nervous people.
\end{failure}

\begin{countermeasures}
  \item Apply it to your own accounts before applying it to anyone else's. Most people's ordinary inaccuracies are larger than they expect.
  \item Do not use a single mismatch as a verdict. Memory genuinely degrades and genuinely reconstructs; the test reads a pattern across several specifics, not one slip.
  \item If you are being tested, notice that the same ground is being covered twice and answer the second time by reference to the first rather than from fresh recall.
  \item Prefer documents. Where a record exists, consistency testing is unnecessary.
\end{countermeasures}

\sources{B1}
\seesources
\seealso{T-06-01, T-06-03, T-07-02}
